\documentclass[11pt,a4paper]{article}

\usepackage[utf8]{inputenc}
\usepackage[T1]{fontenc}
\usepackage[english]{babel}
\usepackage{amsmath, amssymb, amsthm}
\usepackage{geometry}
\usepackage{hyperref}
\usepackage{booktabs}
\usepackage{array}
\usepackage{mathtools}
\usepackage{lmodern}
\usepackage{graphicx}
\usepackage{float}
\usepackage{subcaption}
\usepackage{multirow}

\geometry{margin=2.5cm}

\theoremstyle{definition}
\newtheorem{definition}{Definition}[section]
\newtheorem{theorem}{Theorem}[section]
\newtheorem{proposition}{Proposition}[section]
\newtheorem{corollary}{Corollary}[section]
\newtheorem{remark}{Remark}[section]
\newtheorem{example}{Example}[section]

\newcommand{\cbrt}[1]{\sqrt[3]{#1}}
\newcommand{\prt}[2]{\sqrt[#1]{#2}}
\newcommand{\eps}{\varepsilon}
\newcommand{\CC}{\mathbb{C}}
\newcommand{\RR}{\mathbb{R}}
\newcommand{\ZZ}{\mathbb{Z}}
\newcommand{\NN}{\mathbb{N}}
\newcommand{\Div}{\mathcal{D}}

\title{\textbf{Pandrosion Iteration in the Complex Plane}\\[0.3em]
  \large Basins of attraction, Steffensen acceleration,
  and the Euler product connection}
\author{Ivan \textsc{Besevic}}
\date{April 2026}

\begin{document}

\maketitle

\begin{abstract}
We extend the generalized Pandrosion iteration --- a geometric fixed-point method for computing $p$-th roots via the polynomial $S_p(s) = 1 + s + \cdots + s^{p-1}$ --- from the real line to the complex plane.  For $x \in \CC \setminus \{0\}$ and $p \geq 2$, the iteration $s_{n+1} = 1 - (x-1)/(x \cdot S_p(s_n))$ admits $p$ distinct fixed points $s_k^* = \alpha_k^{-1}$, where $\alpha_k = x^{1/p} e^{2\pi i k / p}$ are the complex $p$-th roots.  We prove that the contraction ratio $\lambda_{p,x}$ extends analytically to $\CC$ and characterize the \emph{convergence region} $\mathcal{C}_p = \{x \in \CC : |\lambda_{p,x}| < 1\}$, showing that divergence occurs only near the negative real axis.  The basins of attraction of the $p$ fixed points exhibit fractal boundaries reminiscent of Julia sets.  We demonstrate that Steffensen's acceleration preserves quadratic convergence in $\CC$, reaching machine precision in 4--5 steps for generic complex targets.  Finally, we observe that the Riemann zeta function admits the formal representation $\zeta(s) = \prod_p \lim_{m \to \infty} S_m(p^{-s})$, identifying each Euler factor as a ``Pandrosion sum at $m = \infty$.''  On the critical line $\mathrm{Re}(s) = \tfrac{1}{2}$, every ratio $p^{-s}$ satisfies $|p^{-s}| = p^{-1/2} < 1$, placing each factor squarely within the Pandrosion convergence regime.
\end{abstract}

\tableofcontents
\bigskip

%=====================================================================
\section{Introduction}
%=====================================================================

In a companion paper~\cite{besevic2026}, we studied the generalized Pandrosion iteration
\begin{equation}
  \label{eq:pandrosion_real}
  s_{n+1} = 1 - \frac{x-1}{x \cdot S_p(s_n)},
  \qquad
  S_p(s) = \sum_{k=0}^{p-1} s^k,
  \qquad
  v_n = x \cdot s_n^{p-1},
\end{equation}
for real $x > 0$ and integer $p \geq 2$.  That construction, rooted in Pandrosion of Alexandria's geometric method for duplicating the cube~\cite{pappus}, converges linearly to $s^* = x^{-1/p}$ with an explicit contraction ratio $\lambda_{p,x}$ depending only on $\alpha = x^{1/p}$.  Steffensen's acceleration transforms the linear method into a quadratically convergent, derivative-free algorithm.

\medskip

A natural question arises: \emph{does Pandrosion's iteration extend to $\CC$?}  Nothing in the formula~\eqref{eq:pandrosion_real} requires $x$ or $s$ to be real --- the geometric sum $S_p$ and the rational map $h(s) = 1 - (x-1)/(x \cdot S_p(s))$ are well-defined for complex arguments.

\medskip

This extension is not merely formal.  By moving to $\CC$, we gain:
\begin{itemize}
  \item \textbf{All $p$ roots at once}: the $p$ branches $\alpha_k = x^{1/p} e^{2\pi ik/p}$ give $p$ fixed points, each with its own basin of attraction.
  \item \textbf{Fractal geometry}: the basin boundaries exhibit self-similar structure, connecting Pandrosion to the theory of Julia sets.
  \item \textbf{A bridge to number theory}: the infinite geometric sum $S_\infty(s) = (1-s)^{-1}$ appears in each factor of the Euler product for the Riemann zeta function, providing a structural (though not analytical) link between Pandrosion and the distribution of primes.
\end{itemize}

\medskip

\textbf{Outline.}
Section~\ref{sec:complex_iteration} defines the complex Pandrosion iteration and identifies its $p$ fixed points.
Section~\ref{sec:complex_lambda} extends the contraction ratio to $\CC$ and characterizes the convergence region.
Section~\ref{sec:basins} studies the basins of attraction and their fractal boundaries.
Section~\ref{sec:steffensen_complex} demonstrates Steffensen acceleration in $\CC$.
Section~\ref{sec:euler} establishes the formal connection to the Euler product.
Section~\ref{sec:divergence} characterizes the divergence boundary.
Section~\ref{sec:conclusion_complex} concludes.

%=====================================================================
\section{Complex Pandrosion iteration}
\label{sec:complex_iteration}
%=====================================================================

\subsection{Definition}

\begin{definition}[Complex Pandrosion iteration]
\label{def:complex_pandrosion}
Let $x \in \CC \setminus \{0,1\}$ and $p \geq 2$ be an integer.  The \textbf{complex Pandrosion iteration} is the map $h : \CC \to \CC$ defined by
\begin{equation}
  \label{eq:complex_pandrosion}
  h(s) = 1 - \frac{x-1}{x \cdot S_p(s)},
  \qquad
  S_p(s) = \sum_{k=0}^{p-1} s^k = \frac{1 - s^p}{1 - s}.
\end{equation}
The output sequence is $v_n = x \cdot s_n^{p-1}$.
\end{definition}

\begin{remark}
The map $h$ is a rational function of $s$ of degree $p-1$ (after clearing the denominator $S_p(s)$).  It has poles at the $(p-1)$-th roots of unity other than $s = 1$ (where $S_p$ vanishes) and a removable singularity at $s = 1$ (where $S_p(1) = p$).
\end{remark}

\subsection{Fixed points}

\begin{theorem}[Complex fixed points]
\label{thm:complex_fixed_points}
The fixed points of $h$ in $\CC$ are exactly
\begin{equation}
  \label{eq:fixed_points_complex}
  s_k^* = \alpha_k^{-1} = x^{-1/p}\, e^{-2\pi i k/p},
  \qquad k = 0, 1, \ldots, p-1,
\end{equation}
where $\alpha_k = x^{1/p} e^{2\pi ik/p}$ are the $p$ complex $p$-th roots of $x$.  The corresponding outputs are $v_k^* = x \cdot (s_k^*)^{p-1} = \alpha_k$.
\end{theorem}

\begin{proof}
At a fixed point $s^*$, we have $h(s^*) = s^*$, i.e.\
\[
  1 - \frac{x-1}{x \cdot S_p(s^*)} = s^*
  \quad \Longleftrightarrow \quad
  x \cdot S_p(s^*) \cdot (1 - s^*) = x - 1.
\]
Since $S_p(s)(1-s) = 1 - s^p$, this gives $x(1 - (s^*)^p) = x - 1$, hence $(s^*)^p = 1/x$, i.e.\ $s^*$ is a $p$-th root of $x^{-1}$.  There are exactly $p$ such roots in $\CC$, given by~\eqref{eq:fixed_points_complex}.  The output formula follows from $v^* = x \cdot (s^*)^{p-1} = x \cdot (x^{-1/p})^{p-1} \cdot e^{-2\pi ik(p-1)/p} = x^{1/p} e^{2\pi ik/p} = \alpha_k$.
\end{proof}

\begin{remark}[Geometry of fixed points]
The $p$ fixed points $s_k^*$ lie on a circle of radius $|x|^{-1/p}$ centered at the origin, equally spaced by angle $2\pi/p$.  For real $x > 1$, the principal fixed point $s_0^* = x^{-1/p} \in (0,1)$ is the one studied in~\cite{besevic2026}.
\end{remark}

\begin{example}
For $p = 3$, $x = 2+i$:
\begin{align*}
  \alpha_0 &= (2+i)^{1/3} \approx 1.2921 + 0.2013i, \\
  \alpha_1 &= \alpha_0 \, e^{2\pi i/3} \approx -0.8204 + 1.0183i, \\
  \alpha_2 &= \alpha_0 \, e^{4\pi i/3} \approx -0.4717 - 1.2196i.
\end{align*}
The three fixed points $s_k^* = \alpha_k^{-1}$ form a triangle in $\CC$.
\end{example}

%=====================================================================
\section{Complex contraction ratio}
\label{sec:complex_lambda}
%=====================================================================

\subsection{Extension to \texorpdfstring{$\CC$}{C}}

\begin{theorem}[Complex contraction ratio]
\label{thm:complex_lambda}
At the principal fixed point $s_0^* = \alpha^{-1}$ with $\alpha = x^{1/p}$ (principal branch), the contraction ratio is
\begin{equation}
  \label{eq:lambda_complex}
  \lambda_{p,x} = h'(s_0^*) = \frac{(\alpha - 1)\displaystyle\sum_{k=1}^{p-1} k\,\alpha^{p-1-k}}{\displaystyle\sum_{k=0}^{p-1} \alpha^k},
  \qquad \alpha = x^{1/p} \in \CC.
\end{equation}
This is the same formula as in the real case~\cite{besevic2026}, now with $\alpha \in \CC$.  The iteration converges locally to $s_0^*$ whenever $|\lambda_{p,x}| < 1$.
\end{theorem}

\begin{proof}
The derivative $h'(s) = \frac{x-1}{x} \cdot \frac{S_p'(s)}{S_p(s)^2}$ is a rational function of $s$.  Evaluating at $s_0^* = \alpha^{-1}$ by the same algebraic manipulations as in~\cite{besevic2026} (which are valid over $\CC$) yields~\eqref{eq:lambda_complex}.
\end{proof}

\begin{remark}
Since $\lambda_{p,x} \in \CC$ in general, the convergence may be \emph{spiraling} (when $\mathrm{Im}(\lambda) \neq 0$): the iterates wind around the fixed point while contracting.  The convergence rate is governed by $|\lambda_{p,x}|$, not $\lambda_{p,x}$ itself.
\end{remark}

\subsection{Convergence region}

\begin{definition}[Convergence region]
\label{def:convergence_region}
The \textbf{convergence region} of the complex Pandrosion iteration is
\[
  \mathcal{C}_p = \{x \in \CC \setminus \{0\} : |\lambda_{p,x}| < 1\}.
\]
\end{definition}

\begin{table}[h]
\centering
\begin{tabular}{rrccc}
\toprule
$\mathrm{Re}(x)$ & $\mathrm{Im}(x)$ & $|\lambda_{3,x}|$ & $\in \mathcal{C}_3$? \\
\midrule
$0.5$ & $0$ & $0.238$ & \checkmark \\
$1$ & $1$ & $0.283$ & \checkmark \\
$2$ & $0$ & $0.220$ & \checkmark \\
$2$ & $1$ & $0.294$ & \checkmark \\
$3$ & $2$ & $0.427$ & \checkmark \\
$5$ & $0$ & $0.468$ & \checkmark \\
$-1$ & $1$ & $0.874$ & \checkmark \\
$-1$ & $0$ & $1.323$ & $\times$ \\
$-2$ & $0$ & $1.260$ & $\times$ \\
$-2$ & $2$ & $0.877$ & \checkmark \\
\bottomrule
\end{tabular}
\caption{Modulus of the contraction ratio $|\lambda_{3,x}|$ for various $x \in \CC$.  Divergence ($|\lambda| \geq 1$) occurs only near the negative real axis.}
\label{tab:lambda_complex}
\end{table}

\begin{proposition}[Structure of the divergence set]
\label{prop:divergence_set}
For each $p \geq 2$, the divergence set $\Div_p = \CC \setminus \mathcal{C}_p$ is a neighborhood of the negative real ray $(-\infty, 0)$.  More precisely:
\begin{enumerate}
  \item $\Div_p$ is symmetric about the real axis: $x \in \Div_p \Leftrightarrow \bar{x} \in \Div_p$.
  \item $\Div_p$ shrinks as $\mathrm{Im}(x)$ increases: for fixed $\mathrm{Re}(x) < 0$, there exists $b_0 > 0$ such that $|\lambda_{p,x}| < 1$ for $|\mathrm{Im}(x)| > b_0$.
  \item The positive real axis $x > 0$ lies entirely in $\mathcal{C}_p$.
\end{enumerate}
\end{proposition}

\begin{proof}
Part~(1) follows from $\overline{\lambda_{p,x}} = \lambda_{p,\bar{x}}$ (conjugate symmetry of the formula).  Part~(3) was established in~\cite{besevic2026}.

For part~(2): fix $a = \mathrm{Re}(x) < 0$ and write $x = a + ib$.  We must show that $|\lambda_{p,x}| < 1$ for all $|b|$ sufficiently large.  Set $\alpha = x^{1/p}$ (principal branch).  As $|b| \to \infty$, we have $|\alpha| \to \infty$ and $\arg(\alpha) \to \pi/(2p) \in (0, \pi/2)$.  In particular, $|\alpha - 1| / |\alpha| \to 1$ and $\alpha \not\in \RR$.  Since $\lambda_{p,x} = (\alpha - 1) N_p(\alpha) / D_p(\alpha)$ with $D_p(\alpha) = \sum_{k=0}^{p-1} \alpha^k$ and $N_p(\alpha) = \sum_{k=1}^{p-1} k\,\alpha^{p-1-k}$, we have
\[
  |\lambda_{p,x}| = \frac{|\alpha - 1| \cdot |N_p(\alpha)|}{|D_p(\alpha)|}.
\]
For $|\alpha| \gg 1$, $D_p(\alpha) \sim \alpha^{p-1}$ and $N_p(\alpha) \sim \alpha^{p-2}$, so $|\lambda| \sim |\alpha - 1|/|\alpha| \to 1$.  However, the approach is from below: writing $\alpha = r e^{i\theta}$ with $\theta = \arg(x)/(p) \in (0, \pi/p)$, a direct estimate using the triangle inequality on $D_p$ and $N_p$ shows that for $r$ large,
\[
  |\lambda_{p,x}| \leq 1 - \frac{C}{r} + O(r^{-2})
\]
for some constant $C > 0$ depending on $p$ and $\theta$.  In particular, $|\lambda| < 1$ for all $r > r_0(\theta)$, and since $r \to \infty$ as $|b| \to \infty$, this establishes the existence of $b_0$.

More precisely: since $\lambda_{p,x}$ is a continuous function of $x$, and $|\lambda_{p,a}| > 1$ for $x = a < 0$ on the real axis, while $|\lambda| < 1$ for $|b|$ large, the intermediate value theorem guarantees the existence of a critical $b_0 > 0$ at which $|\lambda_{p,a+ib_0}| = 1$.
\end{proof}

%=====================================================================
\section{Basins of attraction}
\label{sec:basins}
%=====================================================================

\subsection{Global dynamics}

For each starting point $s_0 \in \CC$, the Pandrosion orbit $(s_n)_{n \geq 0}$ may converge to one of the $p$ fixed points $s_k^*$, diverge to infinity, or enter a cycle.  The \textbf{basin of attraction} of $s_k^*$ is
\[
  \mathcal{B}_k = \{s_0 \in \CC : s_n \to s_k^* \text{ as } n \to \infty\}.
\]

\begin{theorem}[Principal basin]
\label{thm:principal_basin}
For $x \in \mathcal{C}_p$ with $\mathrm{Re}(x) > 0$, the principal basin $\mathcal{B}_0$ (converging to $s_0^* = x^{-1/p}$) contains the real interval $(0, 1)$.
\end{theorem}

\begin{proof}
For real $s_0 \in (0,1)$ and $x > 0$, this is the global convergence theorem from~\cite{besevic2026}.  For complex $x$ with $\mathrm{Re}(x) > 0$, the real interval $(0,1)$ is mapped by $h$ to a curve in $\CC$ that remains in a neighborhood of $s_0^*$ where $|h'| < 1$, ensuring convergence by the contraction mapping principle.
\end{proof}

\subsection{Fractal boundaries}

\begin{figure}[H]
  \centering
  \includegraphics[width=\textwidth]{pandrosion_complex_basins.png}
  \caption{\textbf{Left:} Basin of attraction of the Pandrosion iteration in $\CC$ for $p=3$, $x=2+i$.  The three fixed points $s_k^*$ are marked with stars.  Color indicates which fixed point the orbit converges to.
  \textbf{Right:} Convergence speed (number of iterations to reach $|s_n - s^*| < 10^{-6}$).  Darker regions converge faster.}
  \label{fig:basins}
\end{figure}

When $x \in \CC \setminus \RR$, the basin boundaries between $\mathcal{B}_0, \mathcal{B}_1, \ldots, \mathcal{B}_{p-1}$ typically form \emph{fractal curves}, analogous to the Julia sets of the rational map $h$.  This is expected: $h$ is a rational map of degree $p-1$, and the theory of complex dynamics~\cite{milnor} guarantees that for generic parameters, the Julia set (which coincides with the basin boundaries) has Hausdorff dimension strictly between $1$ and $2$.

\begin{remark}[Comparison with Newton fractals]
Newton's method for $f(z) = z^p - x$ yields the iteration $z_{n+1} = ((p-1)z_n + x/z_n^{p-1})/p$, also a rational map of degree $p$ with $p$ fixed points (the $p$-th roots of $x$).  The Newton fractal for $z^3 - 1$ is one of the most iconic images in complex dynamics.  The Pandrosion fractal is a \emph{new} fractal arising from an entirely different rational map, offering a complementary viewpoint on the geometry of $p$-th root extraction in $\CC$.
\end{remark}

%=====================================================================
\section{Steffensen acceleration in \texorpdfstring{$\CC$}{C}}
\label{sec:steffensen_complex}
%=====================================================================

\subsection{Complex Steffensen--Pandrosion}

\begin{definition}[Complex Steffensen--Pandrosion]
The Steffensen map $T : \CC \to \CC$ is defined by
\begin{equation}
  \label{eq:steffensen_complex}
  T(s) = s - \frac{(h(s) - s)^2}{h(h(s)) - 2h(s) + s},
\end{equation}
where $h$ is the complex Pandrosion map~\eqref{eq:complex_pandrosion}.
\end{definition}

\begin{theorem}[Quadratic convergence in $\CC$]
\label{thm:steffensen_complex}
For any $x \in \mathcal{C}_p$ and any starting point $s_0$ sufficiently close to a fixed point $s_k^*$ with $|\lambda_{p,x}| < 1$ and $\lambda_{p,x} \neq 0$, the Steffensen--Pandrosion iteration converges \emph{quadratically}:
\[
  |T(s) - s_k^*| \leq K_S \cdot |s - s_k^*|^2 + O(|s - s_k^*|^3),
\]
with $K_S = |h''(s_k^*)| / (2|1 - \lambda_{p,x}|)$.
\end{theorem}

\begin{proof}
Steffensen's acceleration is a universal technique for fixed-point iterations~\cite{stoer}.  The standard proof extends to $\CC$ without modification: if $h$ is holomorphic near $s_k^*$ with $h'(s_k^*) = \lambda \neq 0, 1$, then $T'(s_k^*) = 0$, giving quadratic convergence with the stated constant.
\end{proof}

\subsection{Numerical demonstration}

\begin{table}[h]
\centering
\begin{tabular}{ccccc}
\toprule
$p$ & $x$ & $\alpha = x^{1/p}$ & $|\eps_0|$ & Steps to $10^{-15}$ \\
\midrule
$3$ & $2$ (real) & $1.2599$ & $2.71 \times 10^{-1}$ & $3$ \\
$3$ & $2+i$ & $1.2921 + 0.2013i$ & $7.94 \times 10^{-1}$ & $4$ \\
$3$ & $i$ & $0.8660 + 0.5000i$ & $9.01 \times 10^{-1}$ & $5$ \\
$3$ & $3+2i$ & $1.5040 + 0.2986i$ & $7.80 \times 10^{-1}$ & $4$ \\
$4$ & $1+3i$ & $1.2690 + 0.4097i$ & $1.14$ & $5$ \\
\bottomrule
\end{tabular}
\caption{Steffensen--Pandrosion convergence for various complex targets.  Machine precision is reached in 3--5 steps in all cases, confirming quadratic convergence in $\CC$.}
\label{tab:steffensen_complex}
\end{table}

\begin{example}[Detailed trace for $x = 2+i$, $p = 3$]
With $\alpha = (2+i)^{1/3} \approx 1.2921 + 0.2013i$ and $s_0 = 0.5$:
\begin{align*}
  n=0&: \quad |v_0 - \alpha| = 7.94 \times 10^{-1}, \\
  n=1&: \quad |v_1 - \alpha| = 3.31 \times 10^{-2}, \\
  n=2&: \quad |v_2 - \alpha| = 2.45 \times 10^{-5}, \\
  n=3&: \quad |v_3 - \alpha| = 1.38 \times 10^{-11}, \\
  n=4&: \quad |v_4 - \alpha| = 2.50 \times 10^{-16}.
\end{align*}
The doubling of decimal digits at each step is the hallmark of quadratic convergence.
\end{example}

%=====================================================================
\section{The Pandrosion--Euler connection}
\label{sec:euler}
%=====================================================================

\subsection{Euler product as an infinite Pandrosion product}

The Riemann zeta function admits the Euler product representation
\begin{equation}
  \label{eq:euler_product}
  \zeta(s) = \prod_{p \text{ prime}} \frac{1}{1 - p^{-s}},
  \qquad \mathrm{Re}(s) > 1.
\end{equation}
Each factor can be written as an infinite geometric sum:
\[
  \frac{1}{1 - p^{-s}} = \sum_{k=0}^{\infty} p^{-ks} = \lim_{m \to \infty} S_m(p^{-s}).
\]

\begin{proposition}[Euler--Pandrosion identity]
\label{prop:euler_pandrosion}
The Riemann zeta function is a product of Pandrosion sums:
\begin{equation}
  \label{eq:euler_pandrosion}
  \zeta(s) = \prod_{p \text{ prime}} \lim_{m \to \infty} S_m(p^{-s}),
  \qquad \mathrm{Re}(s) > 1.
\end{equation}
Each factor $S_\infty(p^{-s})$ is the ``Pandrosion sum at $m = \infty$'' with ratio $r = p^{-s}$.
\end{proposition}

\begin{proof}
This is a direct rephrasing of the Euler product~\eqref{eq:euler_product} using the identity $\sum_{k=0}^{\infty} r^k = S_\infty(r) = (1-r)^{-1}$ for $|r| < 1$.  The convergence of each geometric series requires $|p^{-s}| = p^{-\mathrm{Re}(s)} < 1$, i.e.\ $\mathrm{Re}(s) > 0$.  The absolute convergence of the infinite product requires $\mathrm{Re}(s) > 1$.
\end{proof}

\subsection{Pandrosion ratios on the critical line}

\begin{proposition}[Critical line ratios]
\label{prop:critical_line}
On the critical line $s = \tfrac{1}{2} + it$, the Pandrosion ratio of each prime satisfies
\begin{equation}
  \label{eq:critical_ratio}
  |p^{-s}| = p^{-1/2} < 1 \quad \text{for all primes } p.
\end{equation}
In particular, each individual Euler factor converges:
\[
  S_\infty(p^{-s}) = \frac{1}{1 - p^{-1/2} e^{-it \ln p}}
\]
is well-defined and bounded.
\end{proposition}

\begin{proof}
$|p^{-s}| = |p^{-1/2 - it}| = p^{-1/2} |p^{-it}| = p^{-1/2}$, since $|p^{-it}| = |e^{-it \ln p}| = 1$.
\end{proof}

\begin{table}[h]
\centering
\begin{tabular}{ccc}
\toprule
Prime $p$ & Ratio $r = p^{-s}$ at $s = \tfrac{1}{2} + 14.135i$ & $|r| = p^{-1/2}$ \\
\midrule
$2$ & $-0.6586 + 0.2575i$ & $0.7071$ \\
$3$ & $-0.5681 - 0.1030i$ & $0.5774$ \\
$5$ & $-0.3248 + 0.3074i$ & $0.4472$ \\
$7$ & $-0.2715 - 0.2630i$ & $0.3780$ \\
$11$ & $-0.2375 - 0.1858i$ & $0.3015$ \\
$13$ & $0.0350 + 0.2751i$ & $0.2774$ \\
\bottomrule
\end{tabular}
\caption{Pandrosion ratios $p^{-s}$ on the critical line at $s \approx \tfrac{1}{2} + 14.135i$ (near the first non-trivial zero of $\zeta$).  Each ratio has modulus $p^{-1/2} < 1$, placing every Euler factor within the Pandrosion convergence regime.}
\label{tab:critical_ratios}
\end{table}

\subsection{The structural gap}

\begin{remark}[What this connection does and does not imply]
\label{rem:honest}
The identity~\eqref{eq:euler_pandrosion} provides a \emph{structural} reinterpretation of $\zeta(s)$ in Pandrosion's language: each prime contributes a convergent geometric sum $S_\infty(p^{-s})$, and the zeta function is their infinite product.

However, this observation does \textbf{not} provide a path to the Riemann Hypothesis.  The key difficulty is:
\begin{itemize}
  \item Each individual factor $(1 - p^{-s})^{-1}$ is \emph{never zero} (it is a single nonvanishing geometric series).
  \item The zeros of $\zeta(s)$ emerge only in the \emph{infinite product} --- they arise from the collective behavior of infinitely many nonzero terms, not from any individual factor.
  \item Understanding \emph{why} these zeros align on $\mathrm{Re}(s) = \tfrac{1}{2}$ requires techniques from analytic number theory (functional equation, explicit formulas, moment methods) that lie beyond the scope of Pandrosion's geometric framework.
\end{itemize}
We record this connection for its conceptual elegance and as a potential source of geometric intuition, while acknowledging that it does not resolve any open conjecture.
\end{remark}

%=====================================================================
\section{The divergence boundary}
\label{sec:divergence}
%=====================================================================

\subsection{Characterization for \texorpdfstring{$p = 2$}{p=2}}

For $p = 2$, the contraction ratio has the closed form $\lambda_{2,x} = (\sqrt{x} - 1)/(\sqrt{x} + 1)$, which extends to $\CC$ as
\[
  \lambda_{2,x} = \frac{x^{1/2} - 1}{x^{1/2} + 1},
  \qquad x^{1/2} = |x|^{1/2} e^{i\arg(x)/2}.
\]

\begin{proposition}[Divergence boundary for $p = 2$]
\label{prop:div_p2}
$|\lambda_{2,x}| = 1$ if and only if $\mathrm{Re}(x^{1/2}) = 0$, i.e.\ $\arg(x) = \pm \pi$.  Thus $\Div_2 = (-\infty, 0]$: Pandrosion's method for square roots diverges exactly on the non-positive real axis.
\end{proposition}

\begin{proof}
Write $\alpha = x^{1/2} = a + bi$.  Then $|\lambda| = |(\alpha-1)/(\alpha+1)|$.  By a standard calculation, $|(\alpha-1)/(\alpha+1)| = 1$ iff $|\alpha-1| = |\alpha+1|$ iff $a = \mathrm{Re}(\alpha) = 0$.  For the principal square root, $\mathrm{Re}(\alpha) = 0$ iff $x \in (-\infty, 0]$.
\end{proof}

\subsection{General case}

For general $p$, the divergence boundary $\partial \Div_p$ is a curve in $\CC$ passing through $x = 0$ and asymptotic to the negative real axis.  Numerical computation shows:

\begin{table}[h]
\centering
\begin{tabular}{ccccc}
\toprule
$p$ & $x=-1$ & $x=-1+i$ & $x=-1+2i$ & $x=-2+2i$ \\
\midrule
$2$ & $1.000$ & $0.673$ & $0.588$ & $0.705$ \\
$3$ & $1.323$ & $0.873$ & $0.745$ & $0.877$ \\
$4$ & $1.474$ & $0.967$ & $0.815$ & $0.951$ \\
$5$ & $1.560$ & $1.020$ & $0.854$ & $0.991$ \\
\bottomrule
\end{tabular}
\caption{Values of $|\lambda_{p,x}|$ near the negative real axis.  The divergence region grows with $p$.}
\label{tab:divergence}
\end{table}

\begin{figure}[H]
  \centering
  \includegraphics[width=0.7\textwidth]{pandrosion_lambda_heatmap.png}
  \caption{Heatmap of $|\lambda_{3,x}|$ in the complex plane.  Blue: $|\lambda| \ll 1$ (fast convergence).  Yellow/red: $|\lambda| \geq 1$ (divergence).  The divergence region (red) clings to the negative real axis.}
  \label{fig:lambda_heatmap}
\end{figure}

%=====================================================================
\section{Conclusion}
\label{sec:conclusion_complex}
%=====================================================================

We have extended Pandrosion's geometric iteration from $\RR^+$ to $\CC$, establishing that:

\begin{itemize}
  \item The complex iteration admits \textbf{$p$ fixed points}, corresponding to the $p$ complex $p$-th roots of $x$, equally spaced on a circle of radius $|x|^{-1/p}$.
  \item The contraction ratio $\lambda_{p,x}$ extends analytically to $\CC$, with $|\lambda_{p,x}| < 1$ for all $x$ outside a neighborhood of the negative real axis.
  \item The \textbf{basins of attraction} exhibit fractal boundaries, providing a new family of rational iteration fractals complementary to Newton fractals.
  \item \textbf{Steffensen acceleration} preserves quadratic convergence in $\CC$, reaching machine precision in 3--5 steps for generic complex targets.
  \item The Riemann zeta function admits the formal representation $\zeta(s) = \prod_p S_\infty(p^{-s})$, identifying each Euler factor as a Pandrosion sum.  On the critical line, every ratio $|p^{-s}| = p^{-1/2} < 1$ lies within the convergence regime.
\end{itemize}

\medskip

\textbf{Perspectives.}
Several directions merit further investigation:
\begin{itemize}
  \item \emph{Fractal dimension}: compute the Hausdorff dimension of the Pandrosion Julia set as a function of $p$ and $x$.
  \item \emph{Multi-ratio Pandrosion}: define an iteration on $\CC^N$ that simultaneously handles $N$ independent geometric ratios, approaching the structure of the partial Euler product.
  \item \emph{Pandrosion zeta iteration}: explore whether applying Pandrosion-type corrections to partial Euler products accelerates the numerical computation of $\zeta(s)$.
  \item \emph{$p$-adic extension}: the geometric sum $S_p$ makes sense over $p$-adic fields, suggesting a Pandrosion construction in non-Archimedean geometry.
\end{itemize}

%=====================================================================
\begin{thebibliography}{9}

\bibitem{besevic2026}
I.~Besevic,
``Generalized Pandrosion residuals: residual profile of an iterative geometric construction for $p$-th roots,''
Preprint, 2026.
DOI: \href{https://doi.org/10.5281/zenodo.19598498}{10.5281/zenodo.19598498}.

\bibitem{pappus}
Pappus of Alexandria,
\textit{Collectio}, Book~III (circa 340~AD).
Critical edition: F.~Hultsch, \textit{Pappi Alexandrini Collectionis quae supersunt},
Weidmann, Berlin, 1876--1878.

\bibitem{stoer}
J.~Stoer and R.~Bulirsch,
\textit{Introduction to Numerical Analysis},
3rd~ed., Springer, New York, 2002.

\bibitem{milnor}
J.~Milnor,
\textit{Dynamics in One Complex Variable},
3rd~ed., Annals of Mathematics Studies~160, Princeton University Press, 2006.

\bibitem{steffensen}
J.~F. Steffensen,
``Remarks on iteration,''
\textit{Skandinavisk Aktuarietidskrift}, vol.~16, pp.~64--72, 1933.

\end{thebibliography}

\end{document}
